%% =============================================================================
%% SKELETON - Section 3: Data, Strategy Construction, and Performance
%% =============================================================================

\subsection{Data and Strategy Construction}
\label{sec:data_strategies}

Three strategies are constructed following the dynamic rolling fund
approach of \citet{duarte2007risk}: each trade is modelled as a separate
fund, with entry triggered when the basis exceeds a threshold and exit at
convergence or maturity. Data are sourced from Bloomberg, ICPL, and ICE;
samples begin in 2005 (CDS--Bond), 2008 (iTraxx), and 2014 (BTP~Italia),
all ending May~2025.

The \textbf{BTP~Italia} strategy trades the inflation-linked basis between
BTP~Italia and nominal BTPs, where the inflation-linked leg is converted
to a fixed-rate exposure via asset swap (ASW).

The \textbf{CDS--Bond Basis} strategy implements the negative basis trade
(long bond, buy CDS protection) on investment-grade European corporates.
The tradable universe is restricted to single-name CDS whose reference
entity belongs to the current on-the-run iTraxx Main series; when a name
exits the index, it is no longer traded, as we assume it loses liquidity.
All CDS are centrally cleared, eliminating counterparty credit risk.

The \textbf{CDS Index Skew} strategy exploits the gap between the iTraxx
index spread and the duration-weighted average of its single-name
constituents. Only on-the-run series are traded, across four indices
(Main, SnrFin, SubFin, Xover).

Entry/exit rules and thresholds are detailed in
Appendix~\ref{app:entry_exit}.

\begin{figure}[H]
\centering
\includegraphics[width=\textwidth]{\figpath/basis_time_series.pdf}
\caption{Market-Level Basis Time Series for Fixed Income Arbitrage Strategies.
For BTP~Italia, the basis is the cross-sectional mean of all outstanding
issues with at least six months to maturity. For iTraxx~Combined, it is
the mean of the four on-the-run sub-index skews. For CDS--Bond~Basis,
it is the mean of the $N=20$ most negative single-name bases in the
iTraxx~Main universe on each day (only bonds with negative basis are
included). Dashed lines indicate entry thresholds.}
\label{fig:basis_time_series}
\end{figure}

\subsection{Index Construction and Weighting Schemes}
\label{sec:index_construction}

Returns are aggregated daily using both equal-weighted
(\citealp{duarte2007risk}) and signal-weighted
(\citealp{rebonato2021convexity}) schemes, where SW weights each trade
proportionally to the basis magnitude at entry. SW is adopted as baseline
due to superior risk-adjusted performance
(Appendix~\ref{app:sw_ew}). All returns are net of transaction costs
computed from observed bid-ask spreads at three regime-dependent tiers
(LOW/MID/HIGH, determined by the level of iTraxx Main 5Y;
Appendix~\ref{app:fees}).
\subsection{Performance Analysis}
\label{sec:performance}

Table~\ref{tab:summary_stats_daily} reports summary statistics for the
signal-weighted indices. All three strategies exhibit positive Sharpe
ratios and moderate drawdowns, with the CDS--Bond basis showing the
longest sample and iTraxx the highest annualised return.

\input{\tablepath/summary_statistics_daily_article}

\begin{figure}[H]
\centering
\includegraphics[width=\textwidth]{\figpath/equity_curves_combined.pdf}
\caption{Cumulative Return Indices and Drawdowns of Fixed Income Arbitrage Strategies (Signal-Weighted).}
\label{fig:equity_curves}
\end{figure}

Table~\ref{tab:mppm} reports manipulation-proof performance measures
\citep{goetzmann2007portfolio}. MPPM penalises strategies that inflate
Sharpe ratios via dynamic trading or return smoothing; all three strategies
retain significant risk-adjusted performance under this metric.

\input{\tablepath/MPPM_analysis_monthly_article}

Table~\ref{tab:moreira_muir} reports results for volatility-managed
portfolios \citep{moreira2017volatility}, which scale exposure inversely
to realised variance. Managed returns improve Sharpe ratios for two of
the three strategies, consistent with time-varying risk premia rather
than a static risk premium story.
\input{\tablepath/moreira_muir_performance_monthly_article}